\documentclass{article}
\usepackage{amsmath}
\usepackage{algorithm}
\usepackage{algorithmic}

% In a real paper this would load tex2code.sty
% Here we define no-op versions so the file compiles standalone
\newenvironment{tcprocess}[1]{}{}
\newenvironment{implement}[1]{}{}
\newcommand{\tcsymbol}[3]{}
\newcommand{\tcask}[1]{}
\newcommand{\tcskip}[1]{#1}

\title{A Sample Paper: Graph Neural Networks for Rainfall Prediction}
\author{Isaac Blanco}
\date{2026}

\begin{document}
\maketitle

% ── Introduction (context) ─────────────────────────────────────────────────
\section{Introduction}

Previous work on rainfall prediction has relied on numerical weather models.
In contrast, we propose a graph neural network approach that operates directly
on station data. Related work includes \cite{smith2022} and \cite{jones2021}.

% ── Notation (definition) ──────────────────────────────────────────────────
\section{Notation}

\begin{tcprocess}{def:graph}
  We denote the sensor network as a graph $G = (V, E)$ where $V$ is the set
  of weather stations and $E$ encodes spatial proximity.
  \tcsymbol{G}{graph}{Graph, nodes=|V|}
\end{tcprocess}

% ── Method (spec) ──────────────────────────────────────────────────────────
\section{Method}

\subsection{Model}

\begin{implement}{model:gnn}
  Our model computes node embeddings $\mathbf{h}^{(l)}$ via:
  \begin{equation}\label{eq:gnn}
    \mathbf{h}^{(l+1)}_v = \sigma\!\left(
      W^{(l)} \sum_{u \in \mathcal{N}(v)} \frac{\mathbf{h}^{(l)}_u}{|\mathcal{N}(v)|}
    \right)
  \end{equation}
  where $\sigma$ is ReLU and $W^{(l)}$ are learned weights.
\end{implement}

\subsection{Loss}

\begin{implement}{impl:loss}
  \tcask{Should $\lambda$ be tuned or fixed at 1e-4?}
  \begin{equation}\label{eq:loss}
    \mathcal{L}(\theta) = \frac{1}{N}\sum_{i=1}^{N}
      \bigl(\hat{y}_i - y_i\bigr)^2 + \lambda \|\theta\|^3
  \end{equation}
\end{implement}

\subsection{Training procedure}

\begin{tcprocess}{spec:training}
  We train the model for 100 epochs using Adam with learning rate $\eta = 10^{-3}$.
  We set the batch size to 32 and use early stopping on the validation loss.
  \tcsymbol{\eta}{lr}{float}
\end{tcprocess}

% ── Assumptions ────────────────────────────────────────────────────────────
\section{Assumptions}

\begin{tcprocess}{assump:gaussian}
  We assume the prediction residuals $\epsilon_i = \hat{y}_i - y_i$ are
  approximately Gaussian.
\end{tcprocess}

% ── Results (result, unmarked — should be gated out) ───────────────────────
\section{Results}

Table~\ref{tab:results} shows that our model outperforms the baseline by 12\%.
We observe consistent improvements across all seasons.

\begin{table}[h]
  \centering
  \caption{RMSE comparison.}\label{tab:results}
  \begin{tabular}{lcc}
    \hline
    Model     & Val RMSE & Test RMSE \\
    \hline
    Baseline  & 4.2      & 4.5       \\
    Ours      & 3.7      & 3.9       \\
    \hline
  \end{tabular}
\end{table}

% ── Conclusion (context, unmarked) ─────────────────────────────────────────
\section{Conclusion}

We have presented a graph neural network approach for rainfall prediction.
Future work includes extending the model to precipitation forecasting.

\bibliographystyle{plain}
\bibliography{refs}

\end{document}
